\documentclass[10pt,a4paper]{article}
\usepackage[utf8]{inputenc}
\usepackage[T1]{fontenc}
\usepackage{lmodern}
\usepackage[margin=2.2cm]{geometry}
\usepackage{amsmath}
\usepackage{graphicx}
\usepackage{booktabs}
\usepackage{caption}
\usepackage{float}
\usepackage{hyperref}
\captionsetup{font=small,labelfont=bf}
\input{values}

\title{Charge-state labelling from U-Net line maps\\[2pt]
\large what was done, and what came out}
\author{\small Generated by \texttt{docs/charge\_states/make\_results.py}\\
\small run: \texttt{\vARun}}
\date{}

\begin{document}
\maketitle

\section{The question}
The U-Net turns a few \% of the gate plane, measured along rays, into a
binary map of transition lines.  What can be done with that map?  Here we
ask whether it is good enough to read off the charge state $(n_1,n_2)$ of
every point of the plane, and how much of the plane gets the right state.

\section{What was done}

\subsection*{Where each array comes from}
For each simulated device, QArray sweeps $V_\mathrm{P1}$ and $V_\mathrm{P2}$
on a $100\times100$ grid and writes three arrays:

\begin{center}
\begin{tabular}{p{4.2cm}p{4.3cm}p{6.2cm}}
\toprule
file & what it is & how it is made \\
\midrule
\texttt{charge\_states.npy} & true $(n_1,n_2)$ per pixel &
  lowest-energy occupation from the capacitance model.  \textbf{No sensor
  involved.}  Used \emph{only} to score. \\
\texttt{ground\_truth\_labels.npy} ($Y$) & true transition lines &
  pixels where $(n_1,n_2)$ changes between neighbours \\
\texttt{charge\_sensing\_data.npy} & what an experiment measures &
  charge-sensor current \\
\bottomrule
\end{tabular}
\end{center}

\subsection*{The pipeline}
\begin{center}
\begin{tabular}{rl}
1. & charge-sensor signal measured along a few rays \\
2. & $\downarrow$ U-Net $\to$ line probability $P$ per pixel \\
3. & $\downarrow$ threshold $\to$ binary line map $B = P > t$ \\
4. & $\downarrow$ \texttt{label\_states}$(B)$ $\to$ $(n_1,n_2)$ per pixel \\
5. & $\downarrow$ compare with \texttt{charge\_states.npy} $\to$ accuracy \\
\end{tabular}
\end{center}

\noindent
The threshold $t$ is the one stored in each \texttt{unet.pt}; it is not
re-tuned.  The labelling function sees only $B$.  It never sees the sensor
signal or the true states.  In short, it:
\begin{enumerate}
\item splits the plane into regions separated by the lines;
\item classifies every border between two regions as dot~1, dot~2 or
  interdot from its slope;
\item decides which side of each border holds the extra hole (holes:
  occupation grows towards negative voltage);
\item gives the top-right region $(0,0)$ and counts across the borders.
\end{enumerate}
The method is described step by step in \texttt{charge\_states\_method.pdf}.

\subsection*{Two inputs, one function}
\begin{itemize}
\item \textbf{full CSD}: \texttt{label\_states}$(Y)$.  Because $Y$ is itself
  derived from \texttt{charge\_states.npy}, this row only checks that the
  labelling function is free of errors.  It is \emph{not} a physics result.
\item \textbf{dense, no ML}: the conventional route.  Every pixel of the
  charge-sensor image is measured, and lines are found without a network
  (Sec.~\ref{sec:dense}); then \texttt{label\_states} is applied as above.
\item \textbf{U-Net}: \texttt{label\_states}$(P>t)$, one row for each of the
  \vAN{} measurement budgets.
\end{itemize}

\subsection*{Metrics (per device, then averaged)}
Pixels within 1\,px of a true line are not scored.
\begin{itemize}
\item \textbf{pixel accuracy}: the fraction of scored pixels with the right
  $(n_1,n_2)$, after the best constant shift $(\Delta n_1,\Delta n_2)$ per
  device;
\item \textbf{state recall}: the fraction of the true states whose pixels
  are mostly ($>50\,\%$) labelled right; \textbf{state precision}: the same,
  seen from the predicted regions;
\item \textbf{(1,1) IoU}: the overlap of the predicted and true $(1,1)$
  regions, the state used for a qubit;
\item \textbf{consistency}: the fraction of redundant borders whose step
  agrees with the labels.  It needs no ground truth.
\end{itemize}

\noindent
\textbf{Baseline.}  Labelling \emph{every} pixel $(0,0)$ already gives
\vBaseMean\,$\pm$\,\vBaseStd\,\% pixel accuracy (median \vBaseMed\,\%),
because the $(0,0)$ region is large.  Pixel accuracy must be read against
this number.

\subsection*{The non-ML reference: lines from the dense sensor image}
\label{sec:dense}
This is the route without a network: measure the whole grid and find the
transitions in the sensor image directly.  Each transition is a small
\emph{step} in the sensor current, on top of the sensor's own smooth
background, which rises steeply towards one corner.  Three conventional
detectors were tried:
\begin{itemize}
\item \textbf{Canny} edges, with optional background subtraction;
\item \textbf{gradient threshold}: $|\nabla Z|$ minus its smoothed copy,
  strongest pixels kept, thinned;
\item \textbf{step residual}: along each axis, the pixel-to-pixel difference
  minus its running median (the local background slope), optionally divided
  by that slope; the strongest pixels are kept.
\end{itemize}
Each detector has its own knobs (smoothing, threshold, background width,
thinning, dilation, frame handling).  All \vADnNSet{} settings were scanned
on the 75 validation devices, and the one with the best mean state recall
was applied unchanged to the test devices.  It was the step residual
(\vADnKnob).  Canny and the gradient threshold failed: the steep sensor
background at the frame dominates their gradients.  The step residual is
the same idea as finding steps along a single ray, done in 2-D.

\begin{figure}[H]
\centering
\includegraphics[width=\textwidth]{figures/dense_examples.pdf}
\caption{The non-ML reference on the three test devices of
Fig.~\ref{fig:ex}.  (a) Dense sensor image (every pixel measured, no
noise).  (b) Step residual.  (c) Line map.  (d) Labels from
\texttt{label\_states}.}
\label{fig:dense}
\end{figure}

\section{Results}
\vANDev{} held-out test devices; the device pool is unfiltered, with 4 to
102 charge states per window (Fig.~\ref{fig:var}).

\begin{table}[H]
\centering
\caption{Test split, \vANDev{} devices.  Each cell: mean [median] over
devices.  First row: the ``all (0,0)'' baseline.}
\label{tab:test}
\small
\resizebox{\textwidth}{!}{\input{table_A_test.tex}}
\end{table}

\begin{figure}[H]
\centering
\includegraphics[width=0.92\textwidth]{figures/both_sweeps.pdf}
\caption{Pixel accuracy and state recall against the fraction of the grid
measured, one point per U-Net.  Diamond: dense sensor image with no network
(100\,\% measured).  Dashed: labelling everything $(0,0)$.  Star: full CSD
$Y$, which only checks the method.}
\label{fig:cov}
\end{figure}

\subsection*{Main findings}
\begin{itemize}
\item \textbf{Without a network, from the dense image, the charge states are
  recovered.}  The step-residual line map gives \vADnAcc\,\% pixel accuracy
  (median \vADnAccMed\,\%), \vADnRec\,\% state recall (median
  \vADnRecMed\,\%), $(1,1)$ IoU \vADnIoU{} and consistency \vADnCon.  This
  needs 100\,\% of the grid measured, and the images are noise-free.
\item \textbf{The U-Net maps, from \vALoCov--\vAHiCov\,\% of the grid, are far
  behind}: \vAHiAcc\,\% pixel accuracy and \vAHiRec\,\% state recall at best
  (\vAHiBud).  In this simulation, measuring everything and finding steps
  conventionally is better for charge-state labelling.  What the U-Net
  saves is measurement time (about 20--60$\times$ fewer points), not
  accuracy.
\item \textbf{The labelling step works on a complete line map.}  On $Y$ it
  gets \vAFullAcc\,\% of the pixels and \vAFullRec\,\% of the states right.
  The limit is the line map, not the labelling.
\item \textbf{From the U-Net maps most charge states are lost.}  State
  recall is \vALoRec\,\% at \vALoCov\,\% coverage (\vALoBud) and \vAHiRec\,\%
  (median \vAHiRecMed\,\%) at \vAHiCov\,\% (\vAHiBud, the best budget of
  this sweep).
\item \textbf{Pixel accuracy is only modestly above the baseline.}  It is
  \vALoAcc\,\% (\vALoBud) to \vAHiAcc\,\% (\vAHiBud), against
  \vBaseMean\,\% for ``all (0,0)''.  For \vAFigBud, \vABeat{} of \vANDev{}
  devices beat their own ``all (0,0)'' score.  Per device it ranges from
  \vAAccMin\,\% to \vAAccMax\,\%.
\item \textbf{The qubit state is usually not found.}  The $(1,1)$ IoU for
  \vAHiBud{} has mean \vAHiIoU{} and median \vAHiIoUMed.
\item \textbf{More measurement helps only a little in this range}
  (1.6--4.7\,\% of the grid, Fig.~\ref{fig:cov}).
\item \textbf{Why it fails:} gaps in the predicted lines.  A gap joins two
  states into one region, and every label beyond it is shifted
  (Fig.~\ref{fig:ex}).  The consistency of the U-Net maps is \vAHiCon{},
  against 1.00 for $Y$, so the failure can be detected without ground truth.
\end{itemize}

\begin{figure}[H]
\centering
\includegraphics[width=\textwidth]{figures/A_fig2_labelling_examples.pdf}
\caption{Good, typical and bad test device for the \vAFigBud{} U-Net.
(a) truth, (b) U-Net line map, (c) labels from (b), (d) errors.}
\label{fig:ex}
\end{figure}

\begin{figure}[H]
\centering
\includegraphics[width=\textwidth]{figures/A_fig4_failures.pdf}
\caption{Left: label errors $(\Delta n_1,\Delta n_2)$ of the wrong pixels,
\vAFigBud{} U-Net, all test devices.  Middle, right: per-device accuracy
against the number of states in the window and the interdot capacitance.}
\label{fig:fail}
\end{figure}

\subsection*{Is the problem too many honeycombs?}
The capacitance range ($d_1g_1, d_2g_2 \in [0.8, 6]$) puts 4 to 102 charge
states in the window, and a few-electron experiment would show far fewer.
Table~\ref{tab:density} splits the test devices by the number of states in
the window.  The U-Nets were trained on the whole range and were not
retrained here.

\begin{table}[H]
\centering
\caption{Test devices grouped by the number of charge states in the window:
mean pixel accuracy and state recall.  ``all (0,0)'' is the baseline on the
same devices.}
\label{tab:density}
\small
\resizebox{\textwidth}{!}{\input{table_A_density.tex}}
\end{table}

\noindent
On sparse windows the U-Net's pixel accuracy is higher, but so is the
baseline, because $(0,0)$ covers more of the plane.  The U-Net's gain over
the baseline is \emph{smallest} on the sparse windows.  Its state recall is
about the same in every group.  So, for these trained U-Nets, the number of
honeycombs is not what limits charge-state labelling: the broken lines are.
Whether a U-Net \emph{trained} only on few-electron windows does better can
only be answered by retraining it (Sec.~\ref{sec:next}).

\begin{figure}[H]
\centering
\includegraphics[width=\textwidth]{figures/fig1_device_variation.pdf}
\caption{How much the test devices differ.}
\label{fig:var}
\end{figure}

\begin{table}[H]
\centering
\caption{Validation split (75 devices).  These devices chose each U-Net's
threshold, so they are \emph{not} held out.}
\small
\resizebox{\textwidth}{!}{\input{table_A_val.tex}}
\end{table}

\section{Caveats and what is missing}
\begin{itemize}
\item \textbf{Noise-free images.}  The dense reference is near-perfect
  partly because the simulated sensor image has no noise.  With
  experimental noise the step residual would degrade, and so would the
  U-Net; how much is not measured here.
\item \textbf{Equal measurement is not compared.}  The dense reference uses
  every pixel and the U-Net 1.6--4.7\,\%.  No non-ML method was run on the
  same sparse ray data.
\item \textbf{The full-CSD row is not an independent result} (see
  Sec.~2).
\item \textbf{The corner assumption}: the top-right region is taken as
  $(0,0)$.  It holds for every device in this simulated pool, not
  necessarily in an experiment.
\item \textbf{Only one sweep.}  The 4--24-ray sweep, which reached 8.9\,\%
  coverage, is no longer in \texttt{results/}, so it is not in this
  document.
\end{itemize}

\section{What could raise the U-Net's accuracy}
\label{sec:next}
These are proposals.  None of them was run, because each one needs new
training (\texttt{run\_0}--\texttt{run\_2}).
\begin{itemize}
\item \textbf{A narrower, experimentally realistic capacitance range}, for
  example a few-electron window with $\le$\,3 holes per dot.  This gives
  fewer, longer lines per window.  Table~\ref{tab:density} does not support
  it as the main fix, but a network trained only on such windows was not
  tested.
\item \textbf{Train for connected lines, not pixels.}  The failure is gaps,
  and the loss (weighted BCE + Dice) does not penalise a one-pixel break.
  A topology-aware loss (e.g.\ clDice) or a skeleton-recall term targets
  exactly this.
\item \textbf{Predict the charge state directly}, as a second U-Net output.
  Then a gap no longer shifts every label behind it.
\item \textbf{More rays rather than more points per ray.}  The earlier
  4--24-ray sweep rose with the number of rays.  Fewer rays will not help.
\item \textbf{Use the consistency as a trigger.}  It flags a broken map
  without ground truth, so a second pass of rays can be aimed where the map
  disagrees with itself.
\end{itemize}

\end{document}
